\pdfoutput=1
\documentclass[11pt]{amsart}

\usepackage[T1]{fontenc}
\usepackage{amsmath,amssymb,amsthm,mathtools}
\usepackage[margin=1.15in]{geometry}
\usepackage{booktabs}
\usepackage{array}
\usepackage{alltt}
\usepackage[hidelinks]{hyperref}
\usepackage{url}
\hypersetup{pdftitle={Bounds for increasing sequences with distinct consecutive sums (Erd\H{o}s Problem 357)},
  pdfauthor={Chandler Snively}}


\newtheorem{theorem}{Theorem}[section]
\newtheorem{proposition}[theorem]{Proposition}
\newtheorem{lemma}[theorem]{Lemma}
\newtheorem{corollary}[theorem]{Corollary}
\theoremstyle{remark}
\newtheorem{remark}[theorem]{Remark}

\setlength{\emergencystretch}{3em}

\newcommand{\Z}{\mathbb{Z}}
\newcommand{\N}{\mathbb{N}}
\newcommand{\isqrt}[1]{\lfloor\sqrt{#1}\rfloor}

\title[Bounds for sequences with distinct consecutive sums]{Bounds for increasing sequences with distinct consecutive sums
(Erd\H{o}s Problem 357)}

\author{Chandler Snively}
\address{Quana Labs}
\email{lab@quanalabs.com}
\urladdr{https://quanalabs.com}

\date{26 October 2026}

\subjclass[2020]{11B13, 11B75, 05D05}
\keywords{distinct consecutive sums, block sums, Sidon sets, Erd\H{o}s problems, formal verification, Lean}

\begin{document}

\begin{abstract}
Let $f(n)$ be the largest $k$ for which there are integers $1\le a_1<\dots<a_k\le n$ such that the sums
$a_i+a_{i+1}+\dots+a_{i'}$ $(1\le i\le i'\le k)$ are pairwise distinct. Whether $f(n)=o(n)$ is open (Erd\H{o}s Problem 357).
We prove
\[
  f(n)\;\le\;\frac{44273}{90000}\,n+2\sqrt n+177792\qquad\text{for all } n\ge0,
\]
so that $\limsup_{n\to\infty} f(n)/n\le\frac{44273}{90000}=\frac12-\frac{2908}{360000}$; the upper bounds on $f$ that we found in
the literature all have leading term at least $n/2$. We also prove
\[
  f(n)\;\ge\;\frac{15}{\sqrt{21}}\sqrt n-1800\qquad\text{for all } n\ge 12096\cdot504^2,
\]
so that $\liminf_{n\to\infty}f(n)/\sqrt n\ge 15/\sqrt{21}=3.27326\ldots$; the best lower bound we found in the literature is
$(4/\sqrt3-o(1))\sqrt n$. The upper bound combines a crossing argument with a lemma that lengthens short blocks by one term, and forces
many integers of $(n/2,2n/3]$ to be missing from the sequence by using pairs of runs of consecutive integers with equal sums. The lower
bound uses a periodic pattern of residues modulo $42$ with one added jump, which reduces every possible coincidence of block sums to a
divisor equation that a counting argument avoids. The paper gives complete written proofs of these bounds. They are also formally
verified in Lean~4 with Mathlib, for a function $f$ that is checked in Lean to be definitionally equal to the one defined in the Formal
Conjectures project; they do not settle the open question whether $f(n)=o(n)$.
\end{abstract}

\maketitle

\section{Introduction}\label{sec:intro}

Call a finite sequence of positive integers $a_1,\dots,a_k$ \emph{valid} if its consecutive sums
\[
  a_i+a_{i+1}+\dots+a_{i'} \qquad (1\le i\le i'\le k)
\]
are pairwise distinct. Following the Erd\H{o}s problems database \cite{EP357}, let
\[
  f(n)=\max\{k:\ \text{there is a valid sequence } 1\le a_1<a_2<\dots<a_k\le n\}.
\]
Problem~357 of \cite{EP357} asks how $f(n)$ grows, and in particular whether $f(n)=o(n)$. Erd\H{o}s and Graham
\cite[p.~58]{ErdosGraham1980} record this question of Erd\H{o}s and Harzheim, together with the variants without monotonicity or without
distinctness; Erd\H{o}s \cite[p.~70]{Erdos1977} discussed the analogous question for infinite sequences. Equivalently, the partial sums $0=s_0<s_1<\dots<s_k$, $s_i=a_1+\dots+a_i$, form a Sidon set whose
consecutive differences increase and are at most $n$. We do not decide whether $f(n)=o(n)$; we prove explicit bounds on both sides.

\begin{theorem}[Upper bound]\label{thm:upper}
For every integer $n\ge0$,
\[
  720000\,f(n)\;\le\;354184\,n+1440000\,\isqrt n+128010240000,
\]
and consequently
\[
  f(n)\;\le\;\frac{44273}{90000}\,n+2\sqrt n+177792,\qquad \frac{44273}{90000}=\frac12-\frac{2908}{360000}=0.4919\overline{2}.
\]
\end{theorem}

Because of its large additive constant, the right-hand side of the second display is below $n/2$ exactly for $n\ge 23202647$. The same method with a
simpler design gives $f(n)\le\frac{797}{1620}n+2\sqrt n+7$ for all $n$ (Remark~\ref{rem:797}; it is proved in Lean, and we give only an outline); its right-hand side is below $n/2$ exactly
for $n\ge63849$, and of the two real forms it gives the smaller bound for every $n\le3348973255$.

\begin{theorem}[Lower bound, modulus $42$]\label{thm:mod42}
For every integer $m\ge0$ we have $f(21m^2+105m+46)\ge 12m$ (for $m=0$ the empty sequence). Consequently, for every $n\ge0$,
\[
  f(n)\;\ge\;\frac{12}{\sqrt{21}}\sqrt n-48,\qquad \frac{12}{\sqrt{21}}=2.618614\ldots .
\]
\end{theorem}

\begin{theorem}[Lower bound, with one skip]\label{thm:skip}
For every integer $j\ge500$ we have $f(12096j^2+87024j+81)\ge 360j$. Consequently
\begin{align*}
  f(n)&\;\ge\;\frac{15}{\sqrt{21}}\sqrt n-1800\qquad\text{for } n\ge 12096\cdot504^2=3{,}072{,}577{,}536,\\
  f(n)&\;\ge\;\frac{15}{\sqrt{21}}\sqrt n-181440\qquad\text{for all } n\ge0,
\end{align*}
where $\frac{15}{\sqrt{21}}=\frac{360}{\sqrt{12096}}=3.273268\ldots$.
\end{theorem}

Here $12096=21\cdot24^2$, so that $360j/\sqrt{12096j^2}=15/\sqrt{21}$. The proofs give the slightly smaller arguments
$21m^2+105m+34$ and $12096j^2+87024j+76$; the forms stated are the ones verified in Lean. The decimal expansions above are truncated,
not rounded; the exact constants are the fractions and radicals.

In particular
\[
  \liminf_{n\to\infty}\frac{f(n)}{\sqrt n}\ge\frac{15}{\sqrt{21}}\qquad\text{and}\qquad\limsup_{n\to\infty}\frac{f(n)}{n}\le\frac{44273}{90000}.
\]
The gap between $\sqrt n$ and $n$ is untouched: nothing in this paper bears on whether $f(n)=o(n)$.

\subsection*{Formal verification}
All statements of Theorems~\ref{thm:upper}--\ref{thm:skip}, and the variant in Remark~\ref{rem:797}, are formally verified in the Lean~4
proof assistant \cite{Lean4} with the Mathlib library \cite{Mathlib2020}, against a definition of $f$ that is definitionally equal to the
one in the Formal Conjectures project \cite{FormalConjectures} (Section~\ref{sec:formal}). The Lean files, the scripts that
repeat every check, and the programs and data of Section~\ref{sec:computations} are in the accompanying public repository \cite{Artifact} (release v1.1). Its release v1.0 (23 October 2026)
contains the first version of this paper; the present version corrects its account of \cite{Hegyvari1986}, corrects some expository
statements found in a further review (among them the hypothesis $t\ge0$ of Proposition~\ref{prop:account}, which its proof uses, and
the form of Theorem~\ref{thm:upper} to which the stated thresholds refer), and shortens the statement on AI use. The theorems, their
proofs and the formal development are unchanged. The written proofs of
Theorems~\ref{thm:upper}--\ref{thm:skip} in Sections~\ref{sec:notation}--\ref{sec:skip} are complete without the formalization;
Remark~\ref{rem:797} is proved only in Lean. The formal proofs follow the same arguments, so they check those arguments mechanically;
they are not independent proofs.

\subsection*{How the results were obtained}
The constructions, the proofs and their formalizations were produced by AI systems, under the direction of the author; the statement on
AI use at the end gives the details, and Section~\ref{sec:discovery} describes how the ideas arose. We report this because it is relevant
to how the results should be checked, and we have arranged the paper so that every step can be checked by hand. The proofs have undergone AI-assisted verification: the
formal verification of Section~\ref{sec:formal}, independent recomputation of every finite step (Section~\ref{sec:computations}), and
adversarial reviews by separate AI systems. They have not received independent human review, and the paper has not been refereed.

\subsection*{Organization}
Section~\ref{sec:previous} reviews previous bounds. Section~\ref{sec:notation} fixes notation and proves two reductions.
Section~\ref{sec:upper} proves Theorem~\ref{thm:upper}; Sections~\ref{sec:mod42} and~\ref{sec:skip} prove Theorems~\ref{thm:mod42}
and~\ref{thm:skip}. Section~\ref{sec:formal} describes the formal verification and what it does and does not certify.
Section~\ref{sec:computations} reports finite computations, which are not used in any proof. Section~\ref{sec:discovery} records how
the results were found, and Section~\ref{sec:limits} discusses limitations.

\section{Previous results}\label{sec:previous}

Table~\ref{tab:lit} lists the bounds on $f$ known to us. Except where marked, we took the bounds from the sources as recorded in
\cite{EP357} and its forum \cite{EP357forum}, without checking them against the original papers. For the non-monotone analogue
$g(n)\ge f(n)$, in which the $a_i$ need not increase, Hegyv\'ari \cite[Theorem~1]{Hegyvari1986} proved
$(\frac13+o(1))n\le g(n)\le(\frac23+o(1))n$ (see also \cite{Beker2024,Konieczny2021}); the upper bound applies to $f$, the lower bound
does not. In the same paper Hegyv\'ari studied shifts of a fixed increasing sequence: for $A=\{1,\dots,k\}$ the least $t$ for which
$\{t+1,\dots,t+k\}$ is valid satisfies $(1+o(1))k^2/25<t<k^2/4$ \cite[Theorem~2]{Hegyvari1986}. The upper estimate, proved with
$t=\lfloor k^2/4\rfloor$, gives $f(n)\ge(2-o(1))\sqrt n$, although Hegyv\'ari does not state it in this form; and the idea of separating
blocks of different lengths by a large common offset, which \cite{Pickhardt2026} and Sections~\ref{sec:mod42}--\ref{sec:skip} combine
with congruence conditions, appears there. Coppersmith and Phillips \cite{CoppersmithPhillips1996} studied the weaker property that no term
is the sum of a block of two or more consecutive earlier terms. A valid sequence has this property, since such a term would collide with the
block consisting of that term alone, so their upper bound applies to $f$; they also constructed sequences with the weaker property of length
$\frac{13}{24}n-O(1)$, so no bound below $\frac{13}{24}n$ can use that property alone (we read only the abstract of
\cite{CoppersmithPhillips1996}).

\begin{table}[ht]
\centering
\small
\renewcommand{\arraystretch}{1.15}
\begin{tabular}{@{}>{\raggedright}p{3.7cm}l>{\raggedright}p{3.9cm}>{\raggedright\arraybackslash}p{3.4cm}@{}}
\toprule
Source & Year & Bound on $f(n)$ & Status\\
\midrule
Hegyv\'ari \cite[Thm.~2]{Hegyvari1986} & 1986 & $\ge(2-o(1))\sqrt n$ & published; implicit in the shifted interval; full text read\\
Straus \cite{Straus1966}, via Weisenberg \cite{EP357forum} & 1966/2025 & $\ge(2+o(1))\sqrt n$ & published construction; observation in a forum comment\\
Pickhardt \cite{Pickhardt2026} & 2026 & $\ge(4/\sqrt3-o(1))\sqrt n$ & web preprint, full text read\\
This paper, Thm.~\ref{thm:mod42} & 2026 & $\ge\frac{12}{\sqrt{21}}\sqrt n-48$ & formally verified\\
This paper, Thm.~\ref{thm:skip} & 2026 & $\ge\frac{15}{\sqrt{21}}\sqrt n-1800$ ($n\ge12096\cdot504^2$) & formally verified\\
\midrule
Hegyv\'ari \cite{Hegyvari1986} (via $f\le g$) & 1986 & $\le(\frac23+o(1))n$ & published (Theorem~1, proved for $g$); full text read\\
Coppersmith--Phillips \cite{CoppersmithPhillips1996}, via Weisenberg \cite{EP357forum} & 1996/2025 & $\le(\frac23-\frac1{512})n+O(\log n)$ & published result; transfer to $f$ in a forum comment\\
Beker \cite{Beker2024}, via Tang \cite{EP357forum} & 2024/2025 & $\le(\sqrt{\tanh 1}+o(1))n$, $\sqrt{\tanh1}=0.87269\ldots$ & published result; deduction in a forum comment; weaker than the two rows above\\
Erd\H{o}s Problem a Day \cite{EPDay2026} & 2026 & $\le \frac n2+O(\sqrt n)$ & informal report; a version of its argument is Corollary~\ref{cor:crossing}\\
Pickhardt \cite{Pickhardt2026}; Lenthall-Cleary \cite{EP357claims} & 2026 & $\le \frac n2+O(n^{2/3})$ & web preprint; unrefereed proof claim\\
This paper, Thm.~\ref{thm:upper} & 2026 & $\le\frac{44273}{90000}n+2\sqrt n+177792$ & formally verified\\
\bottomrule
\end{tabular}
\caption{Bounds on $f(n)$ known to us on 26 October 2026.}
\label{tab:lit}
\end{table}

The lower bound $(2+o(1))\sqrt n$ comes from a stronger property: in an admissible set in the sense of Erd\H{o}s and Straus, sums of
different numbers of distinct elements never coincide, and the largest admissible subset of $\{1,\dots,N\}$ has $(2+o(1))\sqrt N$
elements \cite{Straus1966,DeshouillersFreiman1995,DeshouillersFreiman1999,EP874}. A valid increasing sequence needs the property only
for sums of \emph{consecutive} terms; the constructions of \cite{Pickhardt2026} and of this paper exploit the difference. For small
$n$ the exact values of $f$ are known from computer searches: the least $n$ with $f(n)\ge k$ is tabulated for $k\le22$ in
\cite{OEISA364132}, and $f(n)$ for $n\le53$ in \cite{chy4pro}.

To our knowledge, Theorem~\ref{thm:upper} is the first bound of the form $f(n)\le(c+o(1))n$ with $c<\frac12$; the value $\frac12$
was obtained informally in July 2026 \cite{EPDay2026,Pickhardt2026,EP357claims}. Also to our knowledge, the constants
$\frac{12}{\sqrt{21}}$ and $\frac{15}{\sqrt{21}}$, the constructions of Sections~\ref{sec:mod42}--\ref{sec:skip} and the design of Section~\ref{sec:upper} have not appeared in
the literature; the lower-bound method is Pickhardt's \cite{Pickhardt2026}, with a different pattern and an additional skip. These
statements rest on the literature searches described in the file \path{docs/LITERATURE.md} of the repository \cite{Artifact}, on a
repetition of those searches on 26 October 2026 that found nothing new, and on a reading of the full text of \cite{Hegyvari1986}; such
searches are not a proof of novelty. We could not obtain the full text of \cite{CoppersmithPhillips1996} and read only its abstract.
Theorem~3 of \cite{Hegyvari1986} (see also \cite{EP1213}) shows that an increasing sequence with bounded gaps whose terms are large
enough contains two different blocks with equal sums; Lemma~\ref{lem:pattern}(b) below is the simplest instance of this principle.
Hegyv\'ari does not use it to bound $f$ or $g$; his upper bound on $g$ comes from counting the sums of short blocks. Hegyv\'ari lectured on consecutive sums,
Sidon sets and convex sequences at the CANT conference in 2026; we know only the title and abstract. We cannot exclude unpublished or
unindexed work.

\section{Notation and two reductions}\label{sec:notation}

From now on sequences are indexed from $0$. For a finite sequence $a_0,\dots,a_{k-1}$ and $0\le x\le y\le k$ write
\[
  T(x,y)=a_x+a_{x+1}+\dots+a_{y-1}
\]
for the sum of the \emph{block} $[x,y)$, so $T(x,x)=0$; write $P(y)=T(0,y)$ for the partial sums, and $S(x,\ell)=T(x,x+\ell)$ for the
sum of the block of length $\ell$ starting at $x$. The sequence is valid exactly when $T(x,y)$, for $0\le x<y\le k$, determines
$(x,y)$. A valid strictly increasing sequence in $\{1,\dots,N\}$ is also one in $\{1,\dots,n\}$ for every $n\ge N$, so $f$ is
nondecreasing; and a valid strictly increasing sequence in $\{1,\dots,n\}$ has at most $n$ terms, so the maximum defining $f(n)$
exists.

\begin{lemma}\label{lem:ordered}
Let $0<a_0<a_1<\dots<a_{k-1}$ be integers. If the sequence is not valid, there are $x,\ell,y,s$ with
\[
  1\le s<\ell,\qquad x+\ell\le y,\qquad y+s\le k,\qquad S(x,\ell)=S(y,s),
\]
that is, an earlier, longer block and a later, shorter, disjoint block with the same sum.
\end{lemma}

\begin{proof}
Take two different nonempty blocks with the same sum. Since the blocks are different, have equal sums, and all terms are positive,
neither block contains the other. If they overlap, delete the common part from both: what remains are two nonempty disjoint blocks with
equal sums. So we may assume the blocks are disjoint; call them $(x,\ell)$ and $(y,s)$ with $x+\ell\le y$ (and $y+s\le k$). If
$s\ge\ell$, then $S(y,s)\ge\sum_{i<\ell}a_{y+i}>\sum_{i<\ell}a_{x+i}=S(x,\ell)$, because $a_{y+i}>a_{x+i}$; hence $s<\ell$.
\end{proof}

The formal statement of the problem (Section~\ref{sec:formal}) phrases validity as injectivity of $J\mapsto\sum_{i\in J}a_i$ on the
finite sets $J$ of indices that are \emph{order-connected} (if $i<i'<i''$ and $i,i''\in J$ then $i'\in J$). The next lemma records that
this is the same notion.

\begin{lemma}\label{lem:formal}
Let $a_0<\dots<a_{k-1}$ be positive integers. The map $J\mapsto\sum_{i\in J}a_i$ is injective on the order-connected subsets of
$\{0,\dots,k-1\}$ if and only if the sequence is valid.
\end{lemma}

\begin{proof}
The order-connected subsets are the empty set and the blocks $\{x,\dots,y-1\}$ with $0\le x<y\le k$; distinct pairs $(x,y)$ give distinct
sets. The empty set has sum $0$, and every nonempty block has a positive sum, so the empty set never collides with a block. Hence
injectivity on order-connected sets is injectivity of $(x,y)\mapsto T(x,y)$ on $0\le x<y\le k$.
\end{proof}

\section{The upper bound}\label{sec:upper}

Throughout this section $a_0<a_1<\dots<a_{k-1}$ is a valid sequence of integers in $\{1,\dots,n\}$, and $A=\{a_0,\dots,a_{k-1}\}$.
Since the $a_i$ are distinct positive integers, $a_i\ge i+1$ and $P(y)\ge y(y+1)/2$.

\subsection{Canonical blocks and the crossing bound}

\begin{lemma}[Canonical blocks]\label{lem:canon}
Let $t$ be an integer with $0\le 2t\le n$. For $1\le y\le k$ with $P(y)>t$ let $c(y)$ be the largest $x<y$ with $T(x,y)>t$, and
$\sigma(y)=T(c(y),y)$. Then $t<\sigma(y)\le n$, and $y\mapsto\sigma(y)$ is injective.
\end{lemma}

\begin{proof}
Since $T(0,y)=P(y)>t$, the index $c=c(y)$ exists, and $\sigma(y)>t$ by definition. If $c=y-1$ then $\sigma(y)=a_{y-1}\le n$. Otherwise
$c+1<y$, and the maximality of $c$ gives $T(c+1,y)\le t$; since the block $[c+1,y)$ contains $a_{y-1}>a_c$, we get $a_c<T(c+1,y)\le t$,
so $\sigma(y)=a_c+T(c+1,y)<2t\le n$. If $\sigma(y)=\sigma(y')$, validity gives $[c(y),y)=[c(y'),y')$, so $y=y'$.
\end{proof}

\begin{corollary}[Crossing bound]\label{cor:crossing}
For every $n\ge0$, $f(n)\le\lfloor(n+1)/2\rfloor+\isqrt n$.
\end{corollary}

\begin{proof}
Take $t=\lfloor n/2\rfloor$. The $y\in[1,k]$ with $P(y)>t$ are mapped injectively into $(t,n]$ by Lemma~\ref{lem:canon}, so there are at
most $n-t=\lfloor(n+1)/2\rfloor$ of them. If $P(y)\le t$ then $y^2<y(y+1)\le 2P(y)\le 2t\le n$, so $y\le\isqrt n$; hence at most $\isqrt n$
indices $y$ have $P(y)\le t$.
\end{proof}

This is a version of the argument of \cite{EPDay2026}; we claim no credit for it. It is tight in its main term only when the map $\sigma$ of
Lemma~\ref{lem:canon} covers almost all of $(t,n]$, as it does for $A=(n/2,n]\cap\Z$, where every block with sum at most $n$ is a single
term (that set itself is not valid once $n$ is large). Any further gain must therefore come from coincidences of block sums larger than
$n$; we use runs of a few consecutive integers, whose sums are a few times $n$.

\subsection{Lengthening a block by one term}

\begin{lemma}[One step to the left]\label{lem:chain}
Let $\theta$ be an integer with $3\theta\le 2n+4$. If $1\le x$, $x+2\le y\le k$ and $T(x,y)\le\theta$, then $T(x-1,y)\le n$.
\end{lemma}

\begin{proof}
Put $e=a_{x-1}$ and $v=T(x,y)$. The block $[x,y)$ has at least two terms, and $a_x\ge e+1$, $a_{x+1}\ge e+2$, so $v\ge 2e+3$. Suppose
$e+v>n$. Then $e\ge n+1-v$, so $v\ge 2(n+1-v)+3$, that is $3v\ge 2n+5$. But $3v\le3\theta\le 2n+4$, a contradiction.
\end{proof}

\begin{remark}\label{rem:chainsharp}
The condition $3\theta\le 2n+4$ cannot be weakened to $\theta\le\lfloor(2n+6)/3\rfloor$: for $n=3m$ with $m\ge4$ and the valid sequence $m-1,m,m+2$ the
block $\{m,m+2\}$ has sum $2m+2=\lfloor(2n+6)/3\rfloor$, and adding $m-1$ gives $3m+1>n$. We take $\theta=\lfloor(2n+4)/3\rfloor$ below,
and this is why the holes of Proposition~\ref{prop:design} are counted only in $(n/2,2n/3]$.
\end{remark}

\begin{lemma}[{Leaving $(t,\theta]$}]\label{lem:beyond}
Let $3\theta\le 2n+4$, and let $c+2\le y\le k$ with $T(c,y)\le\theta<P(y)$. Then there is $x'<c$ with $\theta<T(x',y)\le n$.
\end{lemma}

\begin{proof}
Since the terms are positive, $x\mapsto T(x,y)$ is decreasing. Let $x^*$ be the least $x$ with $T(x,y)\le\theta$; then $x^*\le c$, and
$x^*\ge1$ because $T(0,y)=P(y)>\theta$. By minimality $T(x^*-1,y)>\theta$, and Lemma~\ref{lem:chain} (with $x^*+2\le c+2\le y$) gives
$T(x^*-1,y)\le n$. Take $x'=x^*-1<c$.
\end{proof}

\subsection{The accounting}

For integers $t\le\theta$ let $H=\{z\in\Z:t<z\le\theta,\ z\notin A\}$ be the set of \emph{holes} of $A$ in $(t,\theta]$.

\begin{proposition}[Accounting]\label{prop:account}
Let $t,\theta$ be integers with $0\le 2t\le n$, $t\le\theta\le n$ and $3\theta\le 2n+4$. Then
\[
  k+|H|\;\le\;(n-t)+\isqrt{2\theta}.
\]
\end{proposition}

\begin{proof}
Let $\mathcal Y=\{y\in[1,k]:P(y)>t\}$, and let $c$ and $\sigma$ be as in Lemma~\ref{lem:canon}; the set $\Sigma=\sigma(\mathcal Y)$ lies in $(t,n]$ and has
$|\mathcal Y|$ elements. Split $H=H_a\cup H_b$ with $H_a=H\cap\Sigma$ and $H_b=H\setminus\Sigma$. Let
\[
  \mathcal Y_1=\{y\in\mathcal Y:\ \sigma(y)\le\theta,\ \sigma(y)\notin A\},\qquad \mathcal Y_2=\{y\in \mathcal Y_1:\ P(y)>\theta\},\qquad \mathcal Y_1'=\mathcal Y_1\setminus \mathcal Y_2.
\]
Every $z\in H_a$ is $\sigma(y)$ for exactly one $y$, and that $y$ lies in $\mathcal Y_1$; so $|H_a|\le|\mathcal Y_1|$.

Let $y\in \mathcal Y_2$. The canonical block $[c(y),y)$ is not a single term, since otherwise $\sigma(y)=a_{y-1}\in A$; so $c(y)+2\le y$. As
$T(c(y),y)=\sigma(y)\le\theta<P(y)$, Lemma~\ref{lem:beyond} gives $x'_y<c(y)$ with $\Psi(y):=T(x'_y,y)\in(\theta,n]$. By validity, $\Psi$
is injective on $\mathcal Y_2$ (equal sums give equal blocks, hence equal right ends). Moreover $\Psi(y)\notin\Sigma$: if $\Psi(y)=\sigma(y_2)$,
validity gives $[x'_y,y)=[c(y_2),y_2)$, so $y_2=y$ and $x'_y=c(y)$, contradicting $x'_y<c(y)$.

The sets $\Sigma$, $H_b$ and $\Psi(\mathcal Y_2)$ are pairwise disjoint subsets of $(t,n]$: $H_b\subseteq(t,\theta]\subseteq(t,n]$ avoids $\Sigma$ by
definition and avoids $\Psi(\mathcal Y_2)\subseteq(\theta,n]$, and $\Psi(\mathcal Y_2)$ avoids $\Sigma$. Hence
\[
  n-t\;\ge\;|\mathcal Y|+|H_b|+|\mathcal Y_2|\;=\;|\mathcal Y|+|H_b|+|\mathcal Y_1|-|\mathcal Y_1'|\;\ge\;|\mathcal Y|+|H|-|\mathcal Y_1'|.
\]
Finally, $[1,k]\setminus\mathcal Y$ and $\mathcal Y_1'$ are disjoint subsets of $\{y\in[1,k]:P(y)\le\theta\}$, and as in Corollary~\ref{cor:crossing} this
set has at most $\isqrt{2\theta}$ elements. Therefore
\[
  k+|H|=|\mathcal Y|+|[1,k]\setminus\mathcal Y|+|H|\le(n-t)+|\mathcal Y_1'|+|[1,k]\setminus\mathcal Y|\le(n-t)+\isqrt{2\theta}.\qedhere
\]
\end{proof}

\subsection{Runs with equal sums}

For integers $\ell\ge1$ and $w\ge0$ let
\[
  Y_{\ell,w}=\{\ell w,\ \ell w+1,\ \dots,\ \ell w+\ell\},\qquad X_{\ell,w}=\{(\ell+1)w+1,\ \dots,\ (\ell+1)w+\ell\},
\]
runs of $\ell+1$ and of $\ell$ consecutive integers. Their sums are equal:
\[
  \textstyle\sum Y_{\ell,w}=(\ell+1)\ell w+\frac{\ell(\ell+1)}2=\sum X_{\ell,w}.
\]
We call the pair $(\ell,w)$ a \emph{pattern} and $Y_{\ell,w}\cup X_{\ell,w}$ its \emph{support}.

\begin{lemma}\label{lem:pattern}
\textup{(a)} If $c,c+1,\dots,c+r-1$ all lie in $A$, they are consecutive terms $a_i,\dots,a_{i+r-1}$ of the sequence.
\textup{(b)} For no pattern $(\ell,w)$ are $Y_{\ell,w}$ and $X_{\ell,w}$ both contained in $A$.
\textup{(c)} If $t<\ell w$ and $(\ell+1)w+\ell\le\theta$, the support of $(\ell,w)$ lies in $(t,\theta]$ and contains a hole.
\end{lemma}

\begin{proof}
(a) If $c$ and $c+1$ lie in $A$, no term lies strictly between them, so they are consecutive terms; induct on $r$. (b) Otherwise, by (a),
$Y_{\ell,w}$ and $X_{\ell,w}$ are blocks of lengths $\ell+1$ and $\ell$ with equal sums, contradicting validity. (c) The least element of
the support is $\ell w>t$ and the largest is $(\ell+1)w+\ell\le\theta$; by (b) some element of the support is not in $A$.
\end{proof}

\subsection{A design with five bands}

Let $j\ge2$ be an integer and consider $n'=360000j$, with
\[
  t=180000j=n'/2,\qquad \theta=240000j+1=\lfloor(2n'+4)/3\rfloor .
\]
A \emph{band} $(\ell;\alpha,\beta)$ is the set of patterns $(\ell,w)$ with $\alpha j<w\le\beta j$. We use the five bands
\begin{gather*}
  (4;\,45000,45687),\qquad(5;\,36550,37499),\qquad(6;\,31250,31729),\\
  (5;\,38075,39999),\qquad(9;\,22222,23999),
\end{gather*}
which contain $(687+949+479+1924+1777)j=5816j$ patterns in all (the two bands with $\ell=5$ have disjoint ranges of $w$).

\begin{proposition}\label{prop:design}
Let $j\ge2$, and let $a_0<\dots<a_{k-1}$ be a valid sequence in $\{1,\dots,360000j\}$. Then the set $H$ of its holes in
$(180000j,\,240000j+1]$ has at least $2908j$ elements.
\end{proposition}

Only the case $j\ge12$ is used below, and only that case is formalized (\texttt{design\_count}, Section~\ref{sec:formal}); for
$2\le j\le11$ the proof that follows is a written proof only.

\begin{proof}
For a band $(\ell;\alpha,\beta)$ the sets $Y_{\ell,w}$ with $\alpha j<w\le\beta j$ lie in the \emph{$Y$-zone}
$[\ell(\alpha j+1),\ \ell\beta j+\ell]$ and the sets $X_{\ell,w}$ in the \emph{$X$-zone} $[(\ell+1)(\alpha j+1)+1,\ (\ell+1)\beta j+\ell]$.
For the five bands these zones are:
\begin{center}
\small
\begin{tabular}{@{}lll@{}}
\toprule
band & $Y$-zone & $X$-zone\\
\midrule
$(4;45000,45687)$ & $[180000j+4,\ 182748j+4]$ & $[225000j+6,\ 228435j+4]$\\
$(5;36550,37499)$ & $[182750j+5,\ 187495j+5]$ & $[219300j+7,\ 224994j+5]$\\
$(6;31250,31729)$ & $[187500j+6,\ 190374j+6]$ & $[218750j+8,\ 222103j+6]$\\
$(5;38075,39999)$ & $[190375j+5,\ 199995j+5]$ & $[228450j+7,\ 239994j+5]$\\
$(9;22222,23999)$ & $[199998j+9,\ 215991j+9]$ & $[222220j+11,\ 239990j+9]$\\
\bottomrule
\end{tabular}
\end{center}
For $j\ge2$ one checks directly:
\begin{enumerate}
\item every zone lies in $(t,\theta]=(180000j,\,240000j+1]$, so by Lemma~\ref{lem:pattern}(c) each of the $5816j$ patterns has a hole
in its support (the conditions $t<\ell w$ and $(\ell+1)w+\ell\le\theta$ hold at the extreme values of $w$, hence for all $w$ in the band);
\item the five $Y$-zones are pairwise disjoint (in the order listed, each ends before the next begins; the third ends at $190374j+6$, before
$190375j+5$, because $j\ge2$), and all of them lie below every $X$-zone ($215991j+9<218750j+8$);
\item no point lies in three $X$-zones. Write $5$ and $5'$ for the first and second band with $\ell=5$. Ordered by left end, the
$X$-zones are those of $6,5,9,4,5'$, and three inequalities separate them: the zone of $6$ ends before that of $9$ begins
($222103j+6<222220j+11$), the zone of $5$ ends before that of $4$ begins ($224994j+5<225000j+6$), and the zone of $4$ ends before that
of $5'$ begins ($228435j+4<228450j+7$). A zone that ends before another begins also ends before every zone that begins later, so the
only pairs of $X$-zones that can meet are $\{6,5\}$, $\{5,9\}$, $\{9,4\}$ and $\{9,5'\}$. These pairs contain no triangle, and three
intervals on a line with a common point meet pairwise.
\end{enumerate}
Now fix an integer $z$ and count the patterns of the five bands whose support contains $z$. If $z\in Y_{\ell,w}$ then
$\ell w\le z\le\ell w+\ell$, so $w\in\{\lfloor z/\ell\rfloor-1,\lfloor z/\ell\rfloor\}$: at most two patterns of a band, and by (2) only one
band's $Y$-zone and no $X$-zone contains $z$. If $z\in X_{\ell,w}$ then $(\ell+1)w<z<(\ell+1)(w+1)$, so $w=\lfloor z/(\ell+1)\rfloor$: at most
one pattern of a band, and by (3) at most two bands. So every $z$ lies in the supports of at most two patterns.

By (1) choose for each pattern $\omega$ a hole $z_\omega\in H$ in its support. Then
\[
  5816j=\sum_{z\in H}\#\{\omega:z_\omega=z\}\le\sum_{z\in H}\#\{\omega: z\in\operatorname{supp}\omega\}\le2|H|.\qedhere
\]
\end{proof}

\subsection{Proof of Theorem~\ref{thm:upper}}

Let $a_0<\dots<a_{k-1}$ be a valid sequence in $\{1,\dots,n\}$ and $r=\isqrt n$. It suffices to show
$720000k\le354184n+1440000r+128010240000$; dividing by $720000$ and using $r\le\sqrt n$ gives the second form, since
$354184/720000=44273/90000$ and $128010240000/720000=177792$.

\emph{Case $n<4320000$.} By Corollary~\ref{cor:crossing}, $720000k\le360000(n+1)+720000r$, and
$360000(n+1)=354184n+5816n+360000<354184n+5816\cdot4320000+360000<354184n+128010240000$.

\emph{Case $n\ge4320000$.} Let $j=\lceil n/360000\rceil\ge12$ and $n'=360000j$, so $n\le n'\le n+359999$. The sequence is valid in
$\{1,\dots,n'\}$. With $t=180000j$ and $\theta=240000j+1$ we have $2t=n'$, $t\le\theta\le n'$ and $3\theta=2n'+3$, so
Propositions~\ref{prop:account} and~\ref{prop:design} give
\[
  k\le(n'-t)-2908j+\isqrt{2\theta}=177092j+\isqrt{480000j+2}.
\]
Next, $480000j+2=\frac43n'+2\le\frac43(n+359999)+2<\frac43(r+1)^2+480001\le(2r+701)^2$, the last step because
$(2r+701)^2-\frac43(r+1)^2\ge491401-\frac43$ for $r\ge0$. Hence $\isqrt{480000j+2}\le2r+700$, and
\begin{align*}
  720000k&\le354184\,n'+1440000r+504000000\\
  &\le354184n+1440000r+354184\cdot359999+504000000 .
\end{align*}
Since $354184\cdot359999+504000000=128009885816\le128010240000$, the theorem follows. \qed

\begin{remark}[A simpler design]\label{rem:797}
The two-band design with $\ell\in\{5,9\}$, with ranges of $w$ chosen as functions of $n$ rather than of a multiple of $360000$, gives
$1620f(n)\le797n+1620\isqrt{2\lfloor(2n+4)/3\rfloor}+8000$ for all $n$, hence $f(n)\le\frac{797}{1620}n+2\sqrt n+7$; this is formally
verified (\texttt{f\_le\_upper2} and \texttt{f\_le\_real2}, Section~\ref{sec:formal}). We give only an outline, not a complete written
proof. The proof is the one above with Proposition~\ref{prop:design} replaced by a count of the patterns $(5,w)$,
$\lfloor t/5\rfloor<w\le\lfloor(\theta-5)/6\rfloor$, and $(9,w)$, $\lfloor(5\lfloor(\theta-5)/6\rfloor+5)/9\rfloor<w\le\lfloor(\theta-9)/10\rfloor$,
where $t=\lfloor n/2\rfloor$ and $\theta=\lfloor(2n+4)/3\rfloor$; again every integer lies in the supports of at most two patterns. The formal
proof checks the eight floor inequalities this needs, and the closing arithmetic, by linear integer arithmetic for $n\ge400$ and by
evaluation for $12\le n<400$; for $n<12$ the bound follows from Corollary~\ref{cor:crossing}. Nothing else in the paper depends on this
remark except the comparison of the two bounds. Since $\frac{797}{1620}-\frac{44273}{90000}=\frac{43}{810000}$ and both real forms
have the term $2\sqrt n$, the bound of Theorem~\ref{thm:upper} is the smaller of the two exactly when $\frac{43}{810000}n>177792-7$, that is
for $n\ge3348973256$.
\end{remark}

\begin{remark}[What limits this argument]\label{rem:upperlimits}
Two features of the proof cap what it can give. First, the holes are counted only in $(t,\theta]$, and Remark~\ref{rem:chainsharp} shows
that Lemma~\ref{lem:chain} fails beyond $\theta$; a two-term canonical block with sum in $(2n/3,3n/4]$ can have a left neighbour that
makes the sum exceed $n$. Second, the hole count uses only the run pairs $(Y_{\ell,w},X_{\ell,w})$ with supports in zones of bounded
overlap. Numerical experiments recorded in the private research repository (Section~\ref{sec:formal}) are not part of the proofs, are not in
the public repository \cite{Artifact}, and were not re-run for this paper. In them, the linear-programming optimum over interval designs was $0.008105\ldots n$ holes, against $\frac{2908}{360000}n=0.00807\ldots n$
here. That model used run pairs of lengths $p$ and $q$ with $q\le24$ and $1<p/q<4/3$, on a grid of $240$ cells. The fractional
hitting-set optimum over all equal-sum run pairs inside $(t,\theta]$, computed only at $n=3000$, was about $30.3$ holes, that is about
$0.0101n$. These computations are exploratory and cannot be checked from the public repository, and the second, at a single small $n$, does
not give an asymptotic constant. The first suggests that interval designs of this kind cannot do much better than the five bands used
here; we do not claim that the method is limited to any particular constant.
\end{remark}

\section{The modulus-42 pattern}\label{sec:mod42}

Let $(r_0,\dots,r_5)=(2,7,11,25,26,34)$ and $g_t=42\lfloor t/6\rfloor+r_{t\bmod6}$ for $t\ge0$, so that
\[
  (g_t)_{t\ge0}=2,7,11,25,26,34,\;44,49,53,67,68,76,\;86,\dots,\qquad g_{t+6}=g_t+42.
\]
The increments $g_{t+1}-g_t$ are $5,4,14,1,8,10$ periodically, so $g$ is strictly increasing. Write
\begin{equation}\label{eq:eps}
  r_c=7c+\varepsilon_c,\qquad (\varepsilon_0,\dots,\varepsilon_5)=(2,0,-3,4,-2,-1),\qquad \sum_c\varepsilon_c=0,
\end{equation}
and extend $\varepsilon$ $6$-periodically. Then $g_t=7t+\varepsilon_t$ and $-3\le\varepsilon_t\le4$. Put
$\pi(t)=\sum_{i<t}\varepsilon_i$; $\pi$ is $6$-periodic with $(\pi(0),\dots,\pi(5))=(0,2,2,-1,3,1)$, so $-1\le\pi\le3$ and
$\varepsilon_t=\pi(t+1)-\pi(t)$.

\begin{lemma}[Residue spectra]\label{lem:spectra}
Let $w_0,w_1,\dots$ be integers with $w_t\equiv g_t\pmod{42}$ for all $t$. If two blocks of $(w_t)$ have the same sum, their lengths are
congruent modulo $12$.
\end{lemma}

\begin{proof}
Let $G(t)=\sum_{i<t}r_{i\bmod6}\bmod 42$. Since $r_0+\dots+r_5=105\equiv21$, we have $G(t+12)=G(t)$, and
\[
  (G(0),\dots,G(11))=(0,2,9,20,3,29,21,23,30,41,24,8).
\]
A block sum $\sum_{t=x}^{x+\ell-1}w_t$ is congruent to $G(x+\ell)-G(x)$, which lies in the \emph{spectrum}
$D_{\ell\bmod12}$, where $D_\lambda=\{G(x+\lambda)-G(x):0\le x<12\}\subseteq\Z/42\Z$. Direct computation gives
\begin{center}
\small
\begin{tabular}{@{}cl@{\qquad}cl@{}}
$D_0$ & $\{0\}$ & $D_6$ & $\{21\}$\\
$D_1$ & $\{2,7,11,25,26,34\}$ & $D_7$ & $\{4,5,13,23,28,32\}$\\
$D_2$ & $\{9,18,36\}$ & $D_8$ & $\{15,30,39\}$\\
$D_3$ & $\{1,20\}$ & $D_9$ & $\{22,41\}$\\
$D_4$ & $\{3,12,27\}$ & $D_{10}$ & $\{6,24,33\}$\\
$D_5$ & $\{10,14,19,29,37,38\}$ & $D_{11}$ & $\{8,16,17,31,35,40\}$
\end{tabular}
\end{center}
These twelve sets are pairwise disjoint, as one reads off the table, and their sizes add up to $1+6+3+2+3+6+1+6+3+2+3+6=42$, so they
partition $\Z/42\Z$. Hence the residue of a block sum determines its length modulo $12$.
\end{proof}

\begin{proof}[Proof of Theorem~\ref{thm:mod42}]
For $m=0$ there is nothing to prove; let $m\ge1$, $k=12m$, $B=21(m^2+m+2)$ (a multiple of $42$, since $m^2+m$ is even) and
$a_t=B+g_t$ for $0\le t<k$. The sequence is strictly increasing, $a_0=B+2\ge1$, and
$a_{k-1}=B+42(2m-1)+34=21m^2+105m+34$.

Suppose it is not valid, and take $x,\ell,y,s$ as in Lemma~\ref{lem:ordered}. By Lemma~\ref{lem:spectra} (with $w_t=a_t$, since $42\mid B$)
we have $\ell\equiv s\pmod{12}$, so $\ell\ge s+12$. Since $g$ is increasing and nonnegative,
\[
  S(x,\ell)=\ell B+\sum_{t=x}^{x+\ell-1}g_t\ \ge\ \ell B+\sum_{t<\ell}g_t\ \ge\ (s+12)B+\sum_{t<s}g_t .
\]
On the other hand $g_{y+t}-g_t=7y+\varepsilon_{y+t}-\varepsilon_t\le7y+7$, so
\[
  S(y,s)=sB+\sum_{t<s}g_{y+t}\ \le\ sB+\sum_{t<s}g_t+s(7y+7).
\]
Equality of the two sums therefore gives $12B\le s(7y+7)\le 7s(12m+1-s)$, using $y\le k-s$. The right side is at most
$\frac74(12m+1)^2=252m^2+42m+\frac74$, while $12B=252m^2+252m+504$. This contradiction shows the sequence is valid, so
$f(21m^2+105m+46)\ge f(21m^2+105m+34)\ge12m$.

For the bound in $n$, let $Q=\isqrt{\lfloor n/21\rfloor}$. Then $n<21(Q+1)^2$, so $\frac{12}{\sqrt{21}}\sqrt n<12(Q+1)$. If
$Q\le2$ the claimed bound is negative. If $Q\ge3$, put $m=Q-3$; then $21m^2+105m+46\le21(m+3)^2=21Q^2\le n$, so
$f(n)\ge12m=12(Q+1)-48>\frac{12}{\sqrt{21}}\sqrt n-48$.
\end{proof}

\begin{remark}\label{rem:score}
Consider a periodic pattern: an offset plus $q\lfloor t/p\rfloor+r_{t\bmod p}$ for residues $r_0<\dots<r_{p-1}$ modulo $q$, and let the
\emph{gap} $\gamma$ be the least $\gamma\ge1$ such that two blocks of lengths differing by $\gamma$ can have congruent sums modulo $q$ (for
the pattern above, $\gamma=12$ by Lemma~\ref{lem:spectra}). Heuristically, following the plan of \cite{Pickhardt2026}, an offset of about
$\frac{q}{4p\gamma}k^2$ suffices for the size argument when the inseparable length differences are the multiples of $\gamma$, and it
yields the constant $2\sqrt{p\gamma/q}$. We call $p\gamma/q$ the \emph{score} of the pattern. Here $p=6$, $q=42$, $\gamma=12$ give score
$12/7$ and constant $2\sqrt{12/7}=12/\sqrt{21}$; Pickhardt's pattern ($q=3$, $p=2$, $\gamma=2$) gives $4/\sqrt3$, and the modulus-$22$
pattern $(2,6,9,12,18,19)$ ($p=\gamma=6$, score $18/11$) gives $12/\sqrt{22}$. For the modulus-$42$ pattern the offset of
Theorem~\ref{thm:mod42} is close to necessary: numerically, offsets $42t$ still produce coincidences up to $0.84$, $0.94$ and $0.96$
times $21(m^2+m+2)$ for $m=20$, $60$ and $100$.
\end{remark}

\section{The skip construction}\label{sec:skip}

Fix an integer $j\ge1$ and put
\[
  k=360j,\qquad b=288j=\tfrac45k .
\]
For an integer $u\ge0$ consider
\begin{equation}\label{eq:skip}
  a_t=42u+g_t+84\cdot[t\ge b]\qquad(0\le t<k),
\end{equation}
where $[t\ge b]$ is $1$ if $t\ge b$ and $0$ otherwise. The sequence is strictly increasing and positive; adding $84=2\cdot42$ from index $b$
on has the same effect as deleting twelve consecutive terms (two periods) of the pure pattern. Write $\nu(z)=\max(z-b,0)$ for the number
of indices $t<z$ with $t\ge b$.

\begin{lemma}[Block sums]\label{lem:window}
For $x\ge0$, $\ell\ge0$ with $x+\ell\le k$,
\[
  2S(x,\ell)=84u\ell+7\ell(2x+\ell-1)+2\bigl(\pi(x+\ell)-\pi(x)\bigr)+168\bigl(\nu(x+\ell)-\nu(x)\bigr).
\]
\end{lemma}

\begin{proof}
By \eqref{eq:eps}, $a_t=42u+7t+\bigl(\pi(t+1)-\pi(t)\bigr)+84[t\ge b]$. Sum over $x\le t<x+\ell$: the second term gives
$\frac72\ell(2x+\ell-1)$, the third telescopes, and the fourth gives $84(\nu(x+\ell)-\nu(x))$.
\end{proof}

\begin{lemma}[Phase filter]\label{lem:phase}
Let $\xi,\eta,\lambda\in\Z/6\Z$ and
$\Delta=\bigl(\pi(\eta+\lambda)-\pi(\eta)\bigr)-\bigl(\pi(\xi+\lambda)-\pi(\xi)\bigr)$. If $\Delta\equiv0\pmod7$ and
$(\eta-\xi)\lambda+\Delta\equiv0\pmod6$, then $\Delta=0$.
\end{lemma}

\begin{proof}
Since $\pi$ is $6$-periodic, it is a function on $\Z/6\Z$, so $\Delta$ is well defined. If $\lambda=0$, both differences vanish. In general
each difference $\pi(\alpha+\lambda)-\pi(\alpha)$ lies in $[-4,4]$, so $|\Delta|\le8$ and, if $\Delta\neq0$, $\Delta=\pm7$; then one
difference is $\pm4$ and the other $\mp3$. Since $\pi$ takes the value $-1$ only at $3$ and the value $3$ only at $4$, a difference $+4$
occurs only for $(\alpha,\lambda)=(3,1)$ and a difference $-4$ only for $(\alpha,\lambda)=(4,5)$. For $\lambda=1$ the differences are
$\varepsilon_\alpha$, which equals $-3$ only at $\alpha=2$; for $\lambda=5$ they are $-\varepsilon_{\alpha-1}$, which equals $3$ only at
$\alpha=3$. This leaves four cases:
\begin{center}
\small
\begin{tabular}{@{}cccc@{}}
\toprule
$\Delta$ & $(\xi,\eta,\lambda)$ & $(\eta-\xi)\lambda\bmod 6$ & required, $\equiv-\Delta\pmod 6$ \\
\midrule
$+7$ & $(2,3,1)$ & $1$ & $5$\\
$+7$ & $(4,3,5)$ & $1$ & $5$\\
$-7$ & $(3,4,5)$ & $5$ & $1$\\
$-7$ & $(3,2,1)$ & $5$ & $1$\\
\bottomrule
\end{tabular}
\end{center}
(the hypothesis requires $(\eta-\xi)\lambda\equiv-\Delta\pmod 6$). None is possible.
\end{proof}

\begin{proposition}[Reduction to a divisor equation]\label{prop:reduction}
Let $j\ge1$ and $u\ge288j^2+12j+1$. If the sequence \eqref{eq:skip} is not valid, then there are integers $L,\delta$ with
\[
  14\le L\le\delta\le b-1\qquad\text{and}\qquad L\delta=72u+66+12b .
\]
\end{proposition}

\begin{proof}
Take $x,\ell,y,s$ as in Lemma~\ref{lem:ordered}. All of $42u$, $42\lfloor t/6\rfloor$ and $84$ vanish modulo $42$, so
Lemma~\ref{lem:spectra} applies and $h=\ell-s$ is a positive multiple of $12$. Put $d=y-x\ge\ell$. Let
\[
  \Delta=\bigl(\pi(y+s)-\pi(y)\bigr)-\bigl(\pi(x+\ell)-\pi(x)\bigr)\in[-8,8],\qquad
  c_\ell=\nu(x+\ell)-\nu(x),\quad c_s=\nu(y+s)-\nu(y),
\]
so $c_\ell$ and $c_s$ count the indices $\ge b$ in the two blocks. Subtracting the identities of Lemma~\ref{lem:window} for the two blocks
and substituting $\ell=s+h$, $y=x+d$ gives
\begin{equation}\label{eq:K}
  84uh=7\bigl(2sd-2hx-2hs-h^2+h\bigr)+2\Delta+168(c_s-c_\ell).
\end{equation}
As $h\ge1$ and $x,s\ge0$, this implies
\begin{equation}\label{eq:bound}
  84uh\le14sd+16+168(c_s-c_\ell).
\end{equation}
We use $1008u\ge 1008(288j^2+12j+1)=290304j^2+12096j+1008$ and distinguish four cases. They exhaust all possibilities: if
$x\ge b$ we are in (ii); otherwise $x<b$, and either $y\ge b$ (case (iii)) or $y<b$, and then $y+s\le b$ (case (i)) or $y+s>b$
(case (iv)).

\emph{(i) $y+s\le b$.} Then $c_s=c_\ell=0$ and $s+d=y+s-x\le b$, so $sd\le b^2/4=20736j^2$, and \eqref{eq:bound} with $h\ge12$ gives
$1008u\le290304j^2+16$, a contradiction.

\emph{(ii) $x\ge b$.} Then $c_\ell=\ell$, $c_s=s$, so $c_s-c_\ell<0$, and $s+d\le k-b=72j$; \eqref{eq:bound} gives
$1008u\le14\cdot(36j)^2+16$, a contradiction.

\emph{(iii) $x<b\le y$.} Then $c_s=s$, $c_s-c_\ell\le s$, $s\le k-y\le72j$ and $d+s\le k$. By \eqref{eq:bound},
$1008u\le14s(k-s)+168s+16=14s(360j+12-s)+16$. The function $s\mapsto s(360j+12-s)$ is increasing for $s\le180j$, so this is at most
$14\cdot72j\,(288j+12)+16=290304j^2+12096j+16$, a contradiction.

\emph{(iv) $y<b<y+s$.} Then $c_\ell=0$ (the long block ends before $y<b$) and $c_s=y+s-b\le s-1$; in particular $s\ge2$. If $h\ge24$,
\eqref{eq:bound} and $s+d\le k$ give $2016u\le14(180j)^2+168\cdot360j+16=453600j^2+60480j+16$, while
$2016u\ge580608j^2+24192j+2016$; a contradiction. Hence $h=12$, and \eqref{eq:K} with $c_s-c_\ell=x+d+s-b$ simplifies to
\begin{equation}\label{eq:E}
  504u+462+84b=7d(s+12)+\Delta .
\end{equation}
Since $7\mid504,462,84$ we get $\Delta\equiv0\pmod7$; since $6\mid504,462,84$ we get $7d(s+12)+\Delta\equiv ds+\Delta\equiv0\pmod 6$.
Here $\pi(x+\ell)=\pi(x+s)$ because $\ell=s+12$, so $\Delta$ depends only on $x$, $y$, $s$ modulo $6$, and $d\equiv y-x$. By
Lemma~\ref{lem:phase} (with $\xi\equiv x$, $\eta\equiv y$, and the same $\lambda\equiv s$ in both differences), $\Delta=0$, and \eqref{eq:E} becomes
$(s+12)\,d=72u+66+12b$. Take $L=s+12$ and $\delta=d$: then $L\ge14$, $\delta=d\ge\ell=L$, and $\delta=y-x\le y\le b-1$.
\end{proof}

\begin{proposition}[Counting]\label{prop:counting}
Let $j\ge230$ and $U=[\,288j^2+12j+1,\ 288j^2+2012j\,]$, an interval of $2000j$ integers. Some $u\in U$ admits no pair $(L,\delta)$ as in
Proposition~\ref{prop:reduction}.
\end{proposition}

Only the case $j\ge500$ is used below, and only that case is formalized (\texttt{exists\_good}); see Remark~\ref{rem:j230}.

\begin{proof}
Call $u\in U$ \emph{bad} if a pair exists, and put $M(u)=72u+66+12b=72u+66+3456j$. For $u\in U$, $P_1\le M(u)\le P_2$ with
$P_1=20736j^2+4320j+138$ and $P_2=20736j^2+148320j+66$; note $P_2-P_1=72(2000j-1)$. Since $u\mapsto M(u)$ is injective, the number of bad
$u$ is at most the number of pairs $(L,\delta)$ with $14\le L\le\delta\le 288j-1$ and $P_1\le L\delta\le P_2$, $L\delta\equiv66\pmod{72}$.

\emph{Range of $L$.} If $L\le72j-1$ then $L\delta<(L+1)(\delta+1)\le72j\cdot288j=20736j^2<P_1$. If $L\ge144j+515$ then
$L\delta\ge L^2\ge(144j+515)^2>P_2$. So $72j\le L\le144j+514$, i.e.\ $L=72\mu+\rho$ with $0\le\rho<72$ and $j\le\mu\le2j+7$.

\emph{Count for fixed $L$.} Let $\kappa(\rho)$ be the number of residues $r$ modulo $72$ with $\rho r\equiv66\pmod{72}$. For fixed $L$
the cofactor $\delta$ lies in one of $\kappa(L\bmod72)$ residue classes modulo $72$ and in the interval $[P_1/L,P_2/L]$, of length
$72(2000j-1)/L$, so each class contributes at most $\lfloor(2000j-1)/L\rfloor+1$ values. Moreover
$\sum_{\rho<72}\kappa(\rho)=\#\{(\rho,r)\in(\Z/72)^2:\rho r=66\}=8\cdot12=96$, by the Chinese remainder theorem: $\rho r\equiv2\pmod8$
has $8$ solutions and $\rho r\equiv3\pmod 9$ has $12$.

\emph{Summation.} Using $L\ge72\mu$,
\[
  \#\{\text{bad }u\}\le\sum_{\mu=j}^{2j+7}\sum_{\rho<72}\kappa(\rho)\Bigl(\frac{2000j}{72\mu}+1\Bigr)
  =\frac{8000}{3}\,j\sum_{\mu=j}^{2j+7}\frac1\mu+96(j+8).
\]
By $\frac1{\mu}\le\log\mu-\log(\mu-1)$ for $\mu>j$ and $\log(1+x)\le x$,
$\sum_{\mu=j}^{2j+7}\frac1\mu\le\frac1j+\log\frac{2j+7}{j}\le\log2+\frac{9}{2j}$. Hence, with $\log2<0.6931471808$,
\[
  \#\{\text{bad }u\}\le 12000+\tfrac{8000}{3}(\log 2)\,j+96j+768<1944.4\,j+12768<2000j
\]
for $j\ge230$.
\end{proof}

\begin{proof}[Proof of Theorem~\ref{thm:skip}]
Let $j\ge500$ and choose $u\in U$ as in Proposition~\ref{prop:counting}. By Proposition~\ref{prop:reduction} the sequence
\eqref{eq:skip} is valid. It has $k=360j$ terms, the smallest is $a_0=42u+2\ge1$, and the largest is
\[
  a_{k-1}=42u+42(60j-1)+34+84\le42(288j^2+2012j)+2520j+76=12096j^2+87024j+76 .
\]
Hence $f(12096j^2+87024j+81)\ge360j$.

For $n\ge12096\cdot504^2$ let $Q=\isqrt{\lfloor n/12096\rfloor}\ge504$ and $j=Q-4\ge500$. Then
$12096j^2+87024j+81\le12096(j+4)^2=12096Q^2\le n$, so $f(n)\ge360j$. Also $n<12096(Q+1)^2$ and $12096=21\cdot24^2$, so
$\frac{15}{\sqrt{21}}\sqrt n<\frac{15}{\sqrt{21}}\cdot24\sqrt{21}\,(Q+1)=360(Q+1)=360j+1800$. For $n<12096\cdot504^2$,
$\frac{15}{\sqrt{21}}\sqrt n<360\cdot504=181440$, so the all-$n$ form is trivial there.
\end{proof}

\begin{remark}\label{rem:j230}
The proof of Proposition~\ref{prop:counting} works for $j\ge230$; the formal proof uses $j\ge500$. Proposition~\ref{prop:reduction}
holds for every $j\ge1$, and a finite computation (Section~\ref{sec:computations}; not formalized) shows that for every $j\le499$ the
interval $U$ also contains an offset with no pair $(L,\delta)$. With it, $f(12096j^2+87024j+81)\ge360j$ for every $j\ge1$, and the
argument above would give $f(n)\ge\frac{15}{\sqrt{21}}\sqrt n-1800$ for \emph{all} $n$. We do not use this, because the finite computation
is not part of the verified development; the threshold $n\ge12096\cdot504^2$ and the constant $181440$ in Theorem~\ref{thm:skip} are
artifacts of the formal proof.
\end{remark}

\begin{remark}[Where $15/\sqrt{21}$ comes from]\label{rem:constant}
With $k$ and $b$ free, the four cases of Proposition~\ref{prop:reduction} need, to leading order, $u>b^2/288$ (case (i)),
$u>(k-b)^2/288$ (case (ii)), $u>b(k-b)/72$ (case (iii), for $b\ge k/2$) and $u>k^2/576$ (case (iv) with $h\ge24$). The largest of
these is smallest at $b=\frac45k$, where (i) and (iii) both give $u>k^2/450=288j^2$. Then the largest term is about
$42u\approx\frac{7}{75}k^2$, and $k/\sqrt{7k^2/75}=\sqrt{75/7}=15/\sqrt{21}$. The counting step costs only a window of $O(j)$ offsets
and so affects only lower-order terms. This is an optimization of our sufficient conditions, not a proof that the family cannot do
better; the mirror placement $b<k/2$ was not examined. In exact computations for $j=2$ and $j=3$, coincidences not of the divisor type of
case (iv) with $h=12$ occur for offsets up to $u=1131$ and $u=2561$, within $4$ per cent of the left ends $1177$ and $2629$ of $U$. But
some offsets below the left end already give valid sequences (the smallest valid offsets are $u=1069$ for $j=2$ and $u=2457$ for $j=3$,
about $0.91$ and $0.93$ times the left end), so the thresholds are not sharp, and we do not know the limit of the family.
\end{remark}

\section{Formal verification}\label{sec:formal}

\subsection{What was verified}
The formal development consists of six Lean~4 \cite{Lean4} files, which import only Mathlib \cite{Mathlib2020}:
\begin{center}
\small
\begin{tabular}{@{}ll@{}}
\toprule
file & content\\
\midrule
\path{P0357.lean} & definitions; Lemmas~\ref{lem:ordered}, \ref{lem:formal}, \ref{lem:spectra}; Theorem~\ref{thm:mod42}; Corollary~\ref{cor:crossing}\\
\path{P0357Skip.lean} & Section~\ref{sec:skip}; Theorem~\ref{thm:skip}\\
\path{P0357Upper.lean} & Lemmas~\ref{lem:canon}, \ref{lem:chain}, \ref{lem:beyond}; Proposition~\ref{prop:account}\\
\path{P0357Upper2.lean} & Lemma~\ref{lem:pattern}; Remark~\ref{rem:797}\\
\path{P0357Upper4Gen.lean} & double counting over bands (the last display of the proof of Proposition~\ref{prop:design})\\
\path{P0357Upper4.lean} & the five bands, their degree bound, Proposition~\ref{prop:design}, Theorem~\ref{thm:upper}\\
\bottomrule
\end{tabular}
\end{center}
Besides what the table lists, \path{P0357.lean} contains further results on Problem~357 that this paper does not use: the modulus-$22$
pattern, a reformulation of validity through additive energy, density statements for infinite sequences, and examples of arithmetic
progressions. The axiom audit covers them too. The definitions are (with mathematical symbols for the Unicode of the source)
\begin{alltt}\small
def HasDistinctSums \{\(\iota\) \(\alpha\) : Type*\} [Preorder \(\iota\)] [AddCommMonoid \(\alpha\)]
    (a : \(\iota\) \(\to\) \(\alpha\)) : Prop :=
  \{J : Finset \(\iota\) | (J : Set \(\iota\)).OrdConnected\}.InjOn (fun J \(\mapsto\) \(\sum\) x \(\in\) J, a x)

noncomputable def f (n : \(\mathbb{N}\)) : \(\mathbb{N}\) :=
  sSup \{k : \(\mathbb{N}\) | \(\exists\) a : Fin k \(\to\) \(\mathbb{Z}\),
    Set.range a \(\subseteq\) Set.Icc 1 n \(\wedge\) StrictMono a \(\wedge\) HasDistinctSums a\}
\end{alltt}
The verified statements include (namespaces omitted; $\surd$ is the real square root and \texttt{n.sqrt} the integer one)
\begin{alltt}\small
theorem f_le_upper4 (n : \(\mathbb{N}\)) :
    720000 * f n \(\le\) 354184 * n + 1440000 * n.sqrt + 128010240000
theorem f_le_real4 (n : \(\mathbb{N}\)) : (f n : \(\mathbb{R}\)) \(\le\) 44273 / 90000 * n + 2 * \(\surd\)n + 177792
theorem f_le_real2 (n : \(\mathbb{N}\)) : (f n : \(\mathbb{R}\)) \(\le\) 797 / 1620 * n + 2 * \(\surd\)n + 7
theorem mod42_lower (m : \(\mathbb{N}\)) : 12 * m \(\le\) f (21 * m^2 + 105 * m + 46)
theorem mod42_lower_real (n : \(\mathbb{N}\)) : 12 / \(\surd\)21 * \(\surd\)n - 48 \(\le\) (f n : \(\mathbb{R}\))
theorem skip_lower (j : \(\mathbb{N}\)) (hj : 500 \(\le\) j) :
    360 * j \(\le\) f (12096 * j^2 + 87024 * j + 81)
theorem skip_lower_real_large (n : \(\mathbb{N}\)) (hn : 12096 * 504^2 \(\le\) n) :
    15 / \(\surd\)21 * \(\surd\)n - 1800 \(\le\) (f n : \(\mathbb{R}\))
theorem skip_lower_real (n : \(\mathbb{N}\)) : 15 / \(\surd\)21 * \(\surd\)n - 181440 \(\le\) (f n : \(\mathbb{R}\))
\end{alltt}
and the intermediate results \texttt{main\_count} (Proposition~\ref{prop:account}, for arbitrary $t,\theta$ with the stated hypotheses),
\texttt{design\_count} (Proposition~\ref{prop:design}, stated there for $j\ge12$), \texttt{collision\_pair} and \texttt{distinct\_skip}
(Proposition~\ref{prop:reduction}), \texttt{exists\_good} (Proposition~\ref{prop:counting}, for $j\ge500$) and
\texttt{f\_le\_half\_add\_sqrt} (Corollary~\ref{cor:crossing}). The file \path{verification/Headline.lean} of the repository \cite{Artifact} prints every statement used here.

\subsection{Interpretation: does the formal theorem state the problem?}
A proof assistant checks that the proofs prove the stated theorems; it does not check that the statements mean what we intend. We
therefore record the correspondence explicitly.
\begin{enumerate}
\item \emph{Same definition as the community statement.} Our \texttt{HasDistinctSums} and \texttt{f} are definitionally equal to the
definitions in \path{FormalConjectures/ErdosProblems/357.lean} of the Formal Conjectures project \cite{FormalConjectures}, at the commit
\texttt{d16e05a} (12 August 2026) used for the development and at the later commit \texttt{4b5e527} (15 October 2026). This is checked
mechanically: a script copies the two upstream definitions verbatim into a separate namespace, and Lean accepts the equations
\texttt{Erdos357.f = Upstream.f} and \texttt{Erdos357.HasDistinctSums a $\leftrightarrow$ Upstream.HasDistinctSums a} by \texttt{rfl}.
\item \emph{Same notion as in the introduction.} By Lemma~\ref{lem:formal}, \texttt{HasDistinctSums} for a strictly increasing positive
sequence is validity. The formal sequences are maps \texttt{Fin k $\to$ $\Z$} with values in $[1,n]$ and \texttt{StrictMono}, that is,
$1\le a_0<\dots<a_{k-1}\le n$.
\item \emph{The supremum is a maximum.} The supremum is taken in $\N$, where an unbounded set has supremum $0$. The set in question
contains $0$ and is bounded by $n$, so its supremum is its maximum; the formal proofs of the lower bounds establish boundedness before
using the supremum, and the upper bounds are proved for every element of the set, so neither direction relies on the convention.
\item \emph{No special case.} The theorems are stated for the function $f$ itself, for all $n$ (Theorems~\ref{thm:upper}
and~\ref{thm:mod42}) and for all $n\ge 12096\cdot504^2$ or all $n$ with the weaker constant (Theorem~\ref{thm:skip}); they are not stated
for a restricted class of sequences. They do not address the open question, Formal Conjectures' \texttt{erdos\_357.parts.i}
($f(n)=o(n)$), which remains unproved.
\end{enumerate}

\subsection{How it was checked}
We rebuilt the six files in a fresh Lean project with Lean \texttt{v4.27.0} and Mathlib at commit \texttt{a3a10db} (tag
\texttt{v4.27.0}); the build log, the exact toolchain and every check below are in the repository \cite{Artifact}, together with a script
that repeats them.
\begin{itemize}
\item \emph{Axioms.} For every declaration of the six files, not only the theorems listed above, we computed the axioms on which it
depends. All depend only on the standard axioms \texttt{propext}, \texttt{Classical.choice} and \texttt{Quot.sound}; none uses
\texttt{sorryAx} (an unfinished proof) or \texttt{Lean.ofReduceBool} (\texttt{native\_decide}, which trusts compiled code), and the files
declare no axioms.
\item \emph{Kernel replay.} Every declaration was re-checked by the Lean kernel from the compiled files with \texttt{lean4checker}
\cite{lean4checker} (tag \texttt{v4.27.0}). This re-checks the stored proof terms without trusting the elaborator, the tactics or the build
cache. It uses the same kernel implementation as the build, so it is a check of the build, not an independent kernel.
\item \emph{Finite case checks inside the proofs.} The spectra of Lemma~\ref{lem:spectra} are checked by a kernel-evaluated decision over
$12^4$ cases, Lemma~\ref{lem:phase} over its $216$ cases, and the residue count $96$ of Proposition~\ref{prop:counting} by kernel
evaluation; the harmonic estimate uses $\log x\le x-1$ and Mathlib's bound $\log2<0.6931471808$.
\item \emph{Generated code.} The file \path{P0357Upper4.lean} (the five bands and the case split behind the degree bound of
Proposition~\ref{prop:design}, $19$ breakpoints) was generated by a Python script from a data file listing the bands. The script is not
trusted: Lean checks its output like any other file, and a wrong design would fail to compile. The script and data file are in the
repository.
\end{itemize}
The development was written for the research repository in which the results were found, and the files are copied from it unchanged.
That repository is not public; comments in the Lean files that name its research notes are not needed to read or check the proofs.
Three further files of that repository about Problem~357 are not needed here and are omitted: one on side questions, one proving a
conditional reduction, and one recording a negative result on extending Lemma~\ref{lem:chain}.

\section{Finite computations}\label{sec:computations}

The written proofs of Theorems~\ref{thm:upper}--\ref{thm:skip} are complete; the computations in this section are not used in them. They serve as consistency checks and as data for
small $n$. Each item says how it was checked; ``exhaustive search'' results are computational evidence, not formal proofs.

\subsection*{Checks of the proofs}
\begin{enumerate}
\item Every numerical claim and finite case analysis in Sections~\ref{sec:upper}--\ref{sec:skip} was re-checked with exact integer and
rational arithmetic in a script independent of the Lean development: the zones and degree bound of Proposition~\ref{prop:design}
(including a direct count of the patterns through every integer for $2\le j\le6$), the arithmetic of the proof of Theorem~\ref{thm:upper},
the spectra and the $216$ cases of Lemmas~\ref{lem:spectra} and~\ref{lem:phase}, the identity of Lemma~\ref{lem:window}, the constants
of the four cases of Proposition~\ref{prop:reduction}, the count $96$, and the union bound of Proposition~\ref{prop:counting}.
An adversarial review recomputed them again with separately written code, and tested the lemmas of Sections~\ref{sec:notation}
and~\ref{sec:upper} on every valid sequence in $\{1,\dots,22\}$.
\item Theorem~\ref{thm:mod42}: all block sums of the sequence are distinct for $1\le m\le60$.
\item Proposition~\ref{prop:reduction}, for $j=1,2,3,4$ and \emph{every} $u\in U$: all block sums were computed. Every $u$ with no pair
$(L,\delta)$ gives a valid sequence, and every coincidence found for the other values of $u$ is of the type in case (iv): $h=12$, the short
block contains $b$, and $(s+12)d=72u+66+12b$.
\item For every $j\le499$ the interval $U$ contains an offset with no pair $(L,\delta)$ (the smallest proportion of such offsets is
$0.638$, at $j=16$). This is the finite check mentioned in Remark~\ref{rem:j230}.
\item The exact number of bad $u$ is $337{,}866$ of $10^6$ for $j=500$, against the bound $1944.4j+12768$ of
Proposition~\ref{prop:counting}: the counting bound overestimates the number of bad offsets by about a factor $3$, although the
proved bound itself leaves only about $2$ per cent of $U$ to spare.
\end{enumerate}

\subsection*{Exact values for small $n$}
Let $M(k)$ be the least $n$ with $f(n)\ge k$. Exhaustive computer search gives
\begin{align*}
  (M(1),\dots,M(14))&=(1,2,4,5,7,10,12,13,15,18,21,24,25,29),\\
  (M(15),\dots,M(28))&=(30,33,36,38,41,47,50,52,56,60,63,67,69,75),
\end{align*}
and $f(76)=f(77)=28$; together these determine $f(n)$ for every $n\le77$. The values for $k\le22$ agree with \cite{OEISA364132} and
those for $n\le53$ with \cite{chy4pro}. Each value $M(k)$ consists of a lower bound, a valid $k$-term sequence with largest term $M(k)$,
which anyone can check from the definition, and an upper bound, the nonexistence of such a sequence with largest term $M(k)-1$, which
rests on the correctness of the search. The values were first obtained by two exhaustive programs with different pruning rules,
which agree on every value; for $k\le27$ their nonexistence claims were also repeated with all pruning disabled, but the remaining ones
($f(71)=\dots=f(74)=27$ and $f(76)=f(77)=28$) used a pruning rule based on previously computed values $M(r)$. For this paper a third
program was written from the definition alone; the earlier programs were not consulted while it was written, and their results were
used only for the final comparison. Its
only pruning rules are that each term must leave room for the later ones and that each later term must differ from every block sum of
the terms before it, and it uses no previously computed value. It reproduces $M(k)$ for every $k\le28$, the lists of all sequences
attaining them ($67$ in all), and $f(n)$ for every $n\le77$, including the nonexistence claims above. Its pruned search agrees with a
complete unpruned enumeration of all valid sequences in $\{1,\dots,52\}$ and with a brute force over all subsets of $\{1,\dots,34\}$.
All of this is computational evidence. For example,
\[
  2,3,4,6,11,25,38,40,41,47,52,54,56,57,58,60,61,62,64,66,67,68,69,70,71,72,73,75
\]
is one of the two valid $28$-term sequences with largest term $75$.

\subsection*{Greedy sequences}
Greedy constructions produce valid sequences much longer than our lower bounds for every $n$ we could reach. A reverse greedy
construction (take $a=n,n-1,\dots,1$ in turn and keep $a$ as the new first term whenever the sequence stays valid) gives a valid sequence of $26813=5.99\ldots\sqrt n$ terms in $\{1,\dots,n\}$, $n=2\cdot10^7$, against $11676$ terms from
Theorem~\ref{thm:mod42}'s construction at the same $n$. The sequence was regenerated deterministically for this paper, and two separately
written programs confirmed that it is strictly increasing and that all $26813\cdot26814/2=359{,}481{,}891$ of its block sums are distinct. We do not know whether $f(n)/\sqrt n$ is bounded, and finite data do not decide it.

\section{How the results were found}\label{sec:discovery}

The work was done by AI agents, as described in the statement on AI use. We summarize what the records show, because it explains the
shape of the constructions.

\emph{Lower bounds.} Pickhardt's construction is a periodic pattern in the sense of Remark~\ref{rem:score}, with $q=3$ and $p=2$. According to the
agents' records, an exhaustive search on 9 September 2026 over all residue sets modulo $q\le22$, ranked by the score of
Remark~\ref{rem:score}, found the modulus-$22$ pattern $(2,6,9,12,18,19)$, together with a formal proof of
$f(n)\ge\frac{12}{\sqrt{22}}\sqrt n-36$; on 23 October 2026 a branch-and-bound search for $23\le q\le60$ found patterns of score above
$18/11$ only for $q=42$: four residue sets with $p=6$ and $\gamma=12$, namely $R$, $-R$, $R+21$ and $-R+21$ modulo $42$, where
$R=\{2,7,11,25,26,34\}$; and the records report no pattern of score above $12/7$ for $61\le q\le78$. Independent re-runs of the
exhaustive search confirm these results for $q\le34$; the claims for $35\le q\le78$ rest on the records alone, and the search programs
are not part of the repository. These results are computational, and none of them is needed for the proofs. The skip arose from a numerical experiment:
adding a multiple of $q$ to all terms beyond some index keeps every residue, and the smallest admissible offset was smallest with the jump
near $0.8k$. The written proof of Section~\ref{sec:skip} and its formalization came afterwards; an intermediate version used a window of
$4j^2$ offsets and gave the constant $360/\sqrt{12264}=3.2507\ldots$.

\emph{Upper bound.} The crossing argument (Corollary~\ref{cor:crossing}) was taken from \cite{EPDay2026}. The left-extension lemma and the
accounting of Proposition~\ref{prop:account} were found first with a single family of run pairs, $\ell=6$, giving
$f(n)\le\frac{125}{252}n+2\sqrt n+4$; replacing it by two families of degree two gave $\frac{797}{1620}$ (Remark~\ref{rem:797}). A
linear program over band designs, rounded to the grid $360000j$, gave the five bands of Proposition~\ref{prop:design}.

\section{Limitations}\label{sec:limits}

\begin{itemize}
\item The main question, whether $f(n)=o(n)$, is untouched. Our upper bound improves the leading constant from $\frac12$ to
$\frac{44273}{90000}$, a gain of $\frac{727}{90000}=0.00807\ldots$; Remark~\ref{rem:upperlimits} records exploratory evidence that interval designs of this type cannot do much better.
\item The upper bound of Theorem~\ref{thm:upper}, in its real form, is below $n/2$ only for $n\ge23202647$, and its additive constant $177792$ is an
artifact of rounding the design to multiples of $360000$; it was not optimized.
\item The lower bounds are of order $\sqrt n$; finite data suggest that $f(n)$ is considerably larger, and we do not know its order of
growth. The analysis of the one-skip family stops at $15/\sqrt{21}$ (Remark~\ref{rem:constant}); other patterns, several skips or other
jump sizes change the case analysis and have not been explored systematically.
\item The threshold $n\ge12096\cdot504^2$ in Theorem~\ref{thm:skip} is an artifact of the formal proof (Remark~\ref{rem:j230}).
\item Novelty is supported only by the searches described in Section~\ref{sec:previous}; we have read the full text of
\cite{Hegyvari1986} but only the abstract of \cite{CoppersmithPhillips1996}, and no expert in the area has been consulted.
\item The proofs have been checked by machine and by AI systems; no independent human mathematician has checked them.
\end{itemize}

\section*{Statement on the use of AI tools}
\noindent This work was carried out with AI systems, and we describe that use in full. The research repository in which it was done was
set up and directed by the author, who chose the problem, commissioned the runs described below and decided what to publish.
\begin{itemize}
\item The modulus-$22$ predecessor of Theorem~\ref{thm:mod42} was found and formally proved on 9 September 2026 by an autonomous coding
agent (OpenAI Codex, model identifier \texttt{gpt-6-astra}).
\item Theorem~\ref{thm:mod42}, Theorem~\ref{thm:skip}, the first two versions of the upper bound (constants $\frac{125}{252}$ and
$\frac{797}{1620}$) and their formal proofs were produced on 23 October 2026 by autonomous AI agents built on Claude models
(Anthropic; Claude Sonnet~5.5 and Claude Opus~5.5).
\item Theorem~\ref{thm:upper} (the five-band design and its formal proof) was produced later the same day by a further Claude agent
(Claude Sonnet~5.5).
\item The audits (independent rebuilds, axiom audits, kernel replay, independent computations, literature searches and adversarial
reviews), the written proofs in their present form and the text of this paper were produced between 23 and 26 October 2026 with Claude
(Anthropic, model Claude Opus~5.5).
\end{itemize}
The project's records attribute the mathematics and the formal proofs to these AI
systems and record no human mathematical contribution.
The results have not been checked by an independent human mathematician or refereed; they are supported by the formal verification of
Section~\ref{sec:formal}. The author takes full responsibility for the content of this paper, including the accuracy of this statement, and for
correcting or withdrawing any result that is found to be wrong.

\section*{Acknowledgements}
We thank the maintainers of the Erd\H{o}s problems database and the contributors to its forum, and the authors of the Formal Conjectures
project, whose formal statement of Problem~357 we used.

\bibliographystyle{amsplain}
\bibliography{refs}

\end{document}
